\section{Three failures worth reporting}
\label{sec:negative}

\subsection{A validation set with nothing in it}

The scaffold splitter assigns whole scaffold groups to train, validation and test so
no scaffold spans two subsets. On FreeSolv it assigned the validation set
\emph{zero} molecules. The acyclic group --- every molecule with no ring shares the
empty scaffold --- is 320 of 642 molecules. It overflows the training quota,
overflows the validation quota, and lands in test; after that every remaining group
still fits in training, so validation is never reached. The split came out
$0.50 / 0.00 / 0.50$ instead of $0.60 / 0.20 / 0.20$, which silently disabled early
stopping and left half the dataset in test.

Nothing raised an error. The benchmark ran, produced plausible RMSEs, and wrote them
to disk. The bug surfaced only because an unrelated experiment asked the splitter for
a calibration set and got an empty array. Oversized groups are now assigned
largest-first, so the group too large for any held-out subset inflates training
rather than test, and a rebalancing pass moves whole groups into an undersized
validation set. FreeSolv now splits $0.56 / 0.20 / 0.24$. Results computed before the
fix are not comparable with results after it, and the MoleculeNet numbers in
Section~\ref{sec:conformal} were regenerated.

\subsection{A generator that lost to its own control}

The graph diffusion model scored \genBaselineValidity{} validity --- and the
untrained marginal baseline, which samples atoms and bonds from the dataset
frequencies with no model at all, scored higher. A trained model losing to its own
control is not a weak result; it is a broken one, and it was traced rather than
written up.

Two bugs, compounding. ``No bond'' is 83.7\% of atom pairs in a molecular graph, so
under an unweighted cross entropy, predicting ``none'' everywhere is close to
optimal; the denoiser had learned exactly that, producing graphs of mean degree
$0.26$ against a training-set mean of \genTrainDegree{}. Separately, the
valence-constrained sampler began reverse diffusion from the all-zero bond matrix
rather than the marginal noise prior, and since the reverse posterior's likelihood
term anchors $x_{t-1}$ to $x_t$, starting empty dragged the entire trajectory back
towards empty.

With a class-balanced bond loss (exponent \genClassBalance{}, bond weight
\genBondWeight{}) and prior-initialised sampling, mean generated degree rises to
\genModelDegree{} and strict validity to \genModelValidity{}, against
\genBaselineValidity{} for the marginal baseline
(Table~\ref{tab:generation}). Sampling the same trained model
\emph{without} the valence constraint gives \genUnconstrainedValidity{} at mean
degree \genUnconstrainedDegree{}, so the constraint is doing real work rather than
masking a weak denoiser.

\begin{table}[t]
\centering\small
\begin{tabular}{lrrr}
\toprule
sampler & valid & largest frag. & degree \\
\midrule
model + valence & \genModelValidity{} & \genModelFragment{} & \genModelDegree{} \\
model, unconstrained & \genUnconstrainedValidity{} & \genUnconstrainedFragment{} & \genUnconstrainedDegree{} \\
marginal baseline & \genBaselineValidity{} & \genBaselineFragment{} & \genBaselineDegree{} \\
\midrule
training set & --- & --- & \genTrainDegree{} \\
\bottomrule
\end{tabular}
\caption{Generation. Two validity conventions are reported because the literature
uses both: a single connected valence-correct molecule, and the largest connected
fragment. Mean degree against the training mean is the diagnostic that exposed the
collapse.}
\label{tab:generation}
\end{table}

What remains wrong is specific: \genDisconnectedShare{} of samples are rejected for
being in two or more pieces, against \genValenceShare{} for valence errors. The model
has no connectivity prior and the training objective supplies none. That is the next
thing to fix, and it is stated here rather than hidden behind the lenient
fragment-validity convention. These hyperparameters were selected on the validity
metric reported above, with no held-out set, so that figure should be read as an
upper bound.

\subsection{An improvement that did not improve}

The standard construction for adaptive conformal intervals
\cite{papadopoulos2011reliable} fits a separate model to the residuals and uses its
prediction as the difficulty estimate. Implemented here --- an inner fold held out of
training, a second Gaussian process regressed on $\log|r|$ --- it is worse than the
constant-width baseline on every dataset tried: conditional-coverage drift of
\confEsolResidSpread{} on ESOL, \confLipoResidSpread{} on Lipo and
\confFreesolvResidSpread{} on FreeSolv, against \confEsolStdSpread{},
\confLipoStdSpread{} and \confFreesolvStdSpread{} for doing nothing at all, with
wider intervals throughout. It is reported, and not shipped.
